% Generated by tools/build_prior_work.py from the claims ledger; do not edit.
\bibitem{pwambrosio2000functions}
Ambrosio, L., Fusco, N., and Pallara, D. (2000). \emph{Functions of Bounded Variation and Free Discontinuity Problems}. Clarendon Press, Oxford. ISBN 0-19-850245-1.

\bibitem{pwaraki2005extension}
Araki, H., and Zsid{\'{o}}, L. (2005). ``Extension of the structure theorem of Borchers and its application to half-sided modular inclusions.'' \emph{Reviews in Mathematical Physics} \textbf{17}, 491--543. \url{https://doi.org/10.1142/S0129055X05002388}

\bibitem{pwbeardon1983geometry}
Beardon, A. F. (1983). \emph{The Geometry of Discrete Groups}. Springer. \url{https://doi.org/10.1007/978-1-4612-1146-4}

\bibitem{pwbennett1973logical}
Bennett, C. H. (1973). ``Logical Reversibility of Computation.'' \emph{IBM Journal of Research and Development} \textbf{17}, 525--532. \url{https://doi.org/10.1147/rd.176.0525}

\bibitem{pwbertsekas1992data}
Bertsekas, D. P., and Gallager, R. G. (1992). \emph{Data Networks}. Prentice Hall. ISBN 0-13-200916-1.

\bibitem{pwbiggs1993algebraic}
Biggs, N. (1993). \emph{Algebraic Graph Theory}. Cambridge University Press. \url{https://doi.org/10.1017/CBO9780511608704}

\bibitem{pwbisio2016special}
Bisio, A., D'Ariano, G. M., and Perinotti, P. (2016). ``Special relativity in a discrete quantum universe.'' \emph{Physical Review A} \textbf{94}, 042120. \url{https://doi.org/10.1103/PhysRevA.94.042120}

\bibitem{pwbisognano1975duality}
Bisognano, J. J., and Wichmann, E. H. (1975). ``On the duality condition for a Hermitian scalar field.'' \emph{Journal of Mathematical Physics} \textbf{16}, 985--1007. \url{https://doi.org/10.1063/1.522605}

\bibitem{pwbisognano1976duality}
Bisognano, J. J., and Wichmann, E. H. (1976). ``On the duality condition for quantum fields.'' \emph{Journal of Mathematical Physics} \textbf{17}, 303--321. \url{https://doi.org/10.1063/1.522898}

\bibitem{pwblackwell1953equivalent}
Blackwell, D. (1953). ``Equivalent comparisons of experiments.'' \emph{Annals of Mathematical Statistics} \textbf{24}, 265--272. \url{https://doi.org/10.1214/aoms/1177729032}

\bibitem{pwblanco2013relative}
Blanco, D. D., Casini, H., Hung, L. Y., and Myers, R. C. (2013). ``Relative entropy and holography.'' \emph{Journal of High Energy Physics} \textbf{2013}, 60. \url{https://doi.org/10.1007/JHEP08(2013)060}

\bibitem{pwbollinger1996optimal}
Bollinger, J. J., Itano, W. M., Wineland, D. J., and Heinzen, D. J. (1996). ``Optimal frequency measurements with maximally correlated states.'' \emph{Physical Review A} \textbf{54}, R4649--R4652. \url{https://doi.org/10.1103/PhysRevA.54.R4649}

\bibitem{pwborchers1992cpt}
Borchers, H. J. (1992). ``The CPT-theorem in two-dimensional theories of local observables.'' \emph{Communications in Mathematical Physics} \textbf{143}, 315--332. \url{https://doi.org/10.1007/BF02099011}

\bibitem{pwbossavit1988whitney}
Bossavit, A. (1988). ``Whitney forms: a class of finite elements for three-dimensional computations in electromagnetism.'' \emph{IEE Proceedings A: Physical Science, Measurement and Instrumentation, Management and Education, Reviews} \textbf{135}, 493--500. \url{https://doi.org/10.1049/ip-a-1.1988.0077}

\bibitem{pwbossavit1998computational}
Bossavit, A. (1998). \emph{Computational Electromagnetism}. Academic Press. \url{https://doi.org/10.1016/B978-0-12-118710-1.X5000-4}

\bibitem{pwbousso2002holographic}
Bousso, R. (2002). ``The holographic principle.'' \emph{Reviews of Modern Physics} \textbf{74}, 825--874. \url{https://doi.org/10.1103/RevModPhys.74.825}

\bibitem{pwbrandao2019finite}
Brand{\~{a}}o, F. G. S. L., and Kastoryano, M. J. (2019). ``Finite Correlation Length Implies Efficient Preparation of Quantum Thermal States.'' \emph{Communications in Mathematical Physics} \textbf{365}, 1--16. \url{https://doi.org/10.1007/s00220-018-3150-8}

\bibitem{pwbratteli1997operator}
Bratteli, O., and Robinson, D. W. (1997). \emph{Operator Algebras and Quantum Statistical Mechanics 2}. Springer. \url{https://doi.org/10.1007/978-3-662-03444-6}

\bibitem{pwbraunstein1994statistical}
Braunstein, S. L., and Caves, C. M. (1994). ``Statistical distance and the geometry of quantum states.'' \emph{Physical Review Letters} \textbf{72}, 3439--3443. \url{https://doi.org/10.1103/PhysRevLett.72.3439}

\bibitem{pwbroy2001specification}
Broy, M., and St{\o}len, K. (2001). \emph{Specification and Development of Interactive Systems}. Springer. \url{https://doi.org/10.1007/978-1-4613-0091-5}

\bibitem{pwcallen1985thermodynamics}
Callen, H. B. (1985). \emph{Thermodynamics and an Introduction to Thermostatistics}. Wiley. ISBN 0-471-86256-8.

\bibitem{pwcasini2011towards}
Casini, H., Huerta, M., and Myers, R. C. (2011). ``Towards a derivation of holographic entanglement entropy.'' \emph{Journal of High Energy Physics} \textbf{2011}, 36. \url{https://doi.org/10.1007/JHEP05(2011)036}

\bibitem{pwchaitin1982register}
Chaitin, G. J. (1982). ``Register allocation and spilling via graph coloring.'' \emph{ACM SIGPLAN Notices} \textbf{17}, 98--101. \url{https://doi.org/10.1145/872726.806984}

\bibitem{pwchristensen2016introduction}
Christensen, O. (2016). \emph{An Introduction to Frames and Riesz Bases}. Birkh{\"{a}}user/Springer (Applied and Numerical Harmonic Analysis). \url{https://doi.org/10.1007/978-3-319-25613-9}

\bibitem{pwchung1997spectral}
Chung, F. R. K. (1997). \emph{Spectral Graph Theory}. American Mathematical Society. \url{https://doi.org/10.1090/cbms/092}

\bibitem{pwcombescure2012coherent}
Combescure, M., and Robert, D. (2012). \emph{Coherent States and Applications in Mathematical Physics}. Springer. \url{https://doi.org/10.1007/978-94-007-0196-0}

\bibitem{pwconnes1994neumann}
Connes, A., and Rovelli, C. (1994). ``Von Neumann algebra automorphisms and time-thermodynamics relation in generally covariant quantum theories.'' \emph{Classical and Quantum Gravity} \textbf{11}, 2899--2917. \url{https://doi.org/10.1088/0264-9381/11/12/007}

\bibitem{pwconway1999sphere}
Conway, J. H., and Sloane, N. J. A. (1999). \emph{Sphere Packings, Lattices and Groups}. Springer New York. \url{https://doi.org/10.1007/978-1-4757-6568-7}

\bibitem{pwcourant1928partiellen}
Courant, R., Friedrichs, K., and Lewy, H. (1928). ``{\"{U}}ber die partiellen Differenzengleichungen der mathematischen Physik.'' \emph{Mathematische Annalen} \textbf{100}, 32--74. \url{https://doi.org/10.1007/BF01448839}

\bibitem{pwcourant1967partial}
Courant, R., Friedrichs, K., and Lewy, H. (1967). ``On the partial difference equations of mathematical physics.'' \emph{IBM Journal of Research and Development} \textbf{11}, 215--234. \url{https://doi.org/10.1147/rd.112.0215}

\bibitem{pwcover2006elements}
Cover, T. M., and Thomas, J. A. (2006). \emph{Elements of Information Theory}. Wiley. \url{https://doi.org/10.1002/047174882X}

\bibitem{pwdavila2002open}
D{\'{a}}vila, J. (2002). ``On an open question about functions of bounded variation.'' \emph{Calculus of Variations and Partial Differential Equations} \textbf{15}, 519--527. \url{https://doi.org/10.1007/s005260100135}

\bibitem{pwdavis1984methods}
Davis, P. J., and Rabinowitz, P. (1984). \emph{Methods of Numerical Integration}. Academic Press. \url{https://doi.org/10.1016/C2013-0-10566-1}

\bibitem{pwderezinski2013mathematics}
Derezi{\'{n}}ski, J., and G{\'{e}}rard, C. (2013). \emph{Mathematics of Quantization and Quantum Fields}. Cambridge University Press. \url{https://doi.org/10.1017/CBO9780511894541}

\bibitem{pwdiestel2017graph}
Diestel, R. (2017). \emph{Graph Theory}. Springer, Berlin (Graduate Texts in Mathematics). \url{https://doi.org/10.1007/978-3-662-53622-3}

\bibitem{pwdobrushin1992wulff}
Dobrushin, R., Koteck{\'{y}}, R., and Shlosman, S. (1992). \emph{Wulff Construction: A Global Shape from Local Interaction}. American Mathematical Society (Translations of Mathematical Monographs 104). \url{https://doi.org/10.1090/mmono/104}

\bibitem{pwdoob1957conditional}
Doob, J. L. (1957). ``Conditional Brownian motion and the boundary limits of harmonic functions.'' \emph{Bulletin de la Soci{\'{e}}t{\'{e}} Math{\'{e}}matique de France} \textbf{85}, 431--458. \url{https://doi.org/10.24033/bsmf.1494}

\bibitem{pwelaydi2005introduction}
Elaydi, S. (2005). \emph{An Introduction to Difference Equations}. Springer. \url{https://doi.org/10.1007/0-387-27602-5}

\bibitem{pwengl1996regularization}
Engl, H. W., Hanke, M., and Neubauer, A. (1996). \emph{Regularization of Inverse Problems}. Kluwer (Springer Netherlands), Dordrecht. \url{https://doi.org/10.1007/978-94-009-1740-8}

\bibitem{pwevans2010partial}
Evans, L. C. (2010). \emph{Partial Differential Equations}. American Mathematical Society. \url{https://doi.org/10.1090/gsm/019}

\bibitem{pweymard2000finite}
Eymard, R., Gallou{\"{e}}t, T., and Herbin, R. (2000). ``Finite volume methods.'' In \emph{Handbook of Numerical Analysis, Vol. VII}, 713--1018. Elsevier. \url{https://doi.org/10.1016/S1570-8659(00)07005-8}

\bibitem{pwfagin1995reasoning}
Fagin, R., Halpern, J. Y., Moses, Y., and Vardi, M. Y. (1995). \emph{Reasoning About Knowledge}. The MIT Press. \url{https://doi.org/10.7551/mitpress/5803.001.0001}

\bibitem{pwfallat2011totally}
Fallat, S. M., and Johnson, C. R. (2011). \emph{Totally Nonnegative Matrices}. Princeton University Press. \url{https://doi.org/10.1515/9781400839018}

\bibitem{pwfawzi2015quantum}
Fawzi, O., and Renner, R. (2015). ``Quantum conditional mutual information and approximate Markov chains.'' \emph{Communications in Mathematical Physics} \textbf{340}, 575--611. \url{https://doi.org/10.1007/s00220-015-2466-x}

\bibitem{pwfolland1989harmonic}
Folland, G. B. (1989). \emph{Harmonic Analysis in Phase Space}. Princeton University Press. \url{https://doi.org/10.1515/9781400882427}

\bibitem{pwfrancis1999conway}
Francis, G. K., and Weeks, J. R. (1999). ``Conway's ZIP proof.'' \emph{American Mathematical Monthly} \textbf{106}, 393--399. \url{https://doi.org/10.1080/00029890.1999.12005061}

\bibitem{pwfredkin1982conservative}
Fredkin, E., and Toffoli, T. (1982). ``Conservative logic.'' \emph{International Journal of Theoretical Physics} \textbf{21}, 219--253. \url{https://doi.org/10.1007/BF01857727}

\bibitem{pwgibbons2007geometry}
Gibbons, G. W., and Solodukhin, S. N. (2007). ``The geometry of small causal diamonds.'' \emph{Physics Letters B} \textbf{649}, 317--324. \url{https://doi.org/10.1016/j.physletb.2007.03.068}

\bibitem{pwglimm1987quantum}
Glimm, J., and Jaffe, A. (1987). \emph{Quantum Physics: A Functional Integral Point of View}. Springer, New York. \url{https://doi.org/10.1007/978-1-4612-4728-9}

\bibitem{pwgodsil2001algebraic}
Godsil, C., and Royle, G. (2001). \emph{Algebraic Graph Theory}. Springer. \url{https://doi.org/10.1007/978-1-4613-0163-9}

\bibitem{pwgrady2010discrete}
Grady, L. J., and Polimeni, J. R. (2010). \emph{Discrete Calculus}. Springer. \url{https://doi.org/10.1007/978-1-84996-290-2}

\bibitem{pwgreenberger1989going}
Greenberger, D. M., Horne, M. A., and Zeilinger, A. (1989). ``Going Beyond Bell's Theorem.'' In \emph{Bell's Theorem, Quantum Theory and Conceptions of the Universe}, 69--72. Kluwer (Springer Netherlands), Dordrecht. \url{https://doi.org/10.1007/978-94-017-0849-4_10}

\bibitem{pwgriffiths2017introduction}
Griffiths, D. J. (2017). \emph{Introduction to Electrodynamics}. Cambridge University Press. \url{https://doi.org/10.1017/9781108333511}

\bibitem{pwguido1998extensions}
Guido, D., Longo, R., and Wiesbrock, H. W. (1998). ``Extensions of conformal nets and superselection structures.'' \emph{Communications in Mathematical Physics} \textbf{192}, 217--244. \url{https://doi.org/10.1007/s002200050297}

\bibitem{pwhaag1964algebraic}
Haag, R., and Kastler, D. (1964). ``An algebraic approach to quantum field theory.'' \emph{Journal of Mathematical Physics} \textbf{5}, 848--861. \url{https://doi.org/10.1063/1.1704187}

\bibitem{pwhaag1967equilibrium}
Haag, R., Hugenholtz, N. M., and Winnink, M. (1967). ``On the equilibrium states in quantum statistical mechanics.'' \emph{Communications in Mathematical Physics} \textbf{5}, 215--236. \url{https://doi.org/10.1007/BF01646342}

\bibitem{pwhairer2003geometric}
Hairer, E., Lubich, C., and Wanner, G. (2003). ``Geometric numerical integration illustrated by the St{\"{o}}rmer--Verlet method.'' \emph{Acta Numerica} \textbf{12}, 399--450. \url{https://doi.org/10.1017/S0962492902000144}

\bibitem{pwhairer2006geometric}
Hairer, E., Lubich, C., and Wanner, G. (2006). \emph{Geometric Numerical Integration: Structure-Preserving Algorithms for Ordinary Differential Equations}. Springer. \url{https://doi.org/10.1007/3-540-30666-8}

\bibitem{pwhalmos1958finite}
Halmos, P. R. (1958). \emph{Finite-Dimensional Vector Spaces}. Van Nostrand. \url{https://doi.org/10.1007/978-1-4612-6387-6}

\bibitem{pwhardy2008introduction}
Hardy, G. H., and Wright, E. M. (2008). \emph{An Introduction to the Theory of Numbers}. Oxford University Press. ISBN 978-0-19-921986-5.

\bibitem{pwhayden2004structure}
Hayden, P., Jozsa, R., Petz, D., and Winter, A. (2004). ``Structure of states which satisfy strong subadditivity of quantum entropy with equality.'' \emph{Communications in Mathematical Physics} \textbf{246}, 359--374. \url{https://doi.org/10.1007/s00220-004-1049-z}

\bibitem{pwheinosaari2011mathematical}
Heinosaari, T., and Ziman, M. (2011). \emph{The Mathematical Language of Quantum Theory: From Uncertainty to Entanglement}. Cambridge University Press. \url{https://doi.org/10.1017/CBO9781139031103}

\bibitem{pwhelgason1999radon}
Helgason, S. (1999). \emph{The Radon Transform}. Birkh{\"{a}}user, Boston (Progress in Mathematics). \url{https://doi.org/10.1007/978-1-4757-1463-0}

\bibitem{pwhigham2002accuracy}
Higham, N. J. (2002). \emph{Accuracy and Stability of Numerical Algorithms}. SIAM. \url{https://doi.org/10.1137/1.9780898718027}

\bibitem{pwhodges1993model}
Hodges, W. (1993). \emph{Model Theory}. Cambridge University Press. \url{https://doi.org/10.1017/CBO9780511551574}

\bibitem{pwholevo2012quantum}
Holevo, A. S. (2012). \emph{Quantum Systems, Channels, Information: A Mathematical Introduction}. De Gruyter, Berlin. \url{https://doi.org/10.1515/9783110273403}

\bibitem{pwhorn2012matrix}
Horn, R. A., and Johnson, C. R. (2012). \emph{Matrix Analysis}. Cambridge University Press. \url{https://doi.org/10.1017/CBO9781139020411}

\bibitem{pwjackson1999classical}
Jackson, J. D. (1999). \emph{Classical Electrodynamics}. Wiley, New York. ISBN 0-471-30932-X.

\bibitem{pwjacobson1995thermodynamics}
Jacobson, T. (1995). ``Thermodynamics of Spacetime: The Einstein Equation of State.'' \emph{Physical Review Letters} \textbf{75}, 1260--1263. \url{https://doi.org/10.1103/PhysRevLett.75.1260}

\bibitem{pwjacobson2016entanglement}
Jacobson, T. (2016). ``Entanglement Equilibrium and the Einstein Equation.'' \emph{Physical Review Letters} \textbf{116}, 201101. \url{https://doi.org/10.1103/PhysRevLett.116.201101}

\bibitem{pwkalman1963mathematical}
Kalman, R. E. (1963). ``Mathematical description of linear dynamical systems.'' \emph{Journal of the Society for Industrial and Applied Mathematics, Series A: Control} \textbf{1}, 152--192. \url{https://doi.org/10.1137/0301010}

\bibitem{pwkato2019quantum}
Kato, K., and Brand{\~{a}}o, F. G. S. L. (2019). ``Quantum Approximate Markov Chains are Thermal.'' \emph{Communications in Mathematical Physics} \textbf{370}, 117--149. \url{https://doi.org/10.1007/s00220-019-03485-6}

\bibitem{pwkelly2011reversibility}
Kelly, F. P. (2011). \emph{Reversibility and Stochastic Networks}. Cambridge University Press. ISBN 978-1-107-40115-0.

\bibitem{pwkemeny1976finite}
Kemeny, J. G., and Snell, J. L. (1976). \emph{Finite Markov Chains}. Springer-Verlag, New York (Undergraduate Texts in Mathematics). ISBN 0-387-90192-2.

\bibitem{pwkolmogoroff1936theorie}
Kolmogoroff, A. (1936). ``Zur Theorie der Markoffschen Ketten.'' \emph{Mathematische Annalen} \textbf{112}, 155--160. \url{https://doi.org/10.1007/BF01565412}

\bibitem{pwkraus1983states}
Kraus, K. (1983). \emph{States, Effects, and Operations: Fundamental Notions of Quantum Theory}. Springer. \url{https://doi.org/10.1007/3-540-12732-1}

\bibitem{pwkuipers1974uniform}
Kuipers, L., and Niederreiter, H. (1974). \emph{Uniform Distribution of Sequences}. Wiley. ISBN 0-486-45019-8.

\bibitem{pwkullback1951information}
Kullback, S., and Leibler, R. A. (1951). ``On information and sufficiency.'' \emph{Annals of Mathematical Statistics} \textbf{22}, 79--86. \url{https://doi.org/10.1214/aoms/1177729694}

\bibitem{pwkuo1992discrete}
Kuo, H. J., and Trudinger, N. S. (1992). ``Discrete methods for fully nonlinear elliptic equations.'' \emph{SIAM Journal on Numerical Analysis} \textbf{29}, 123--135. \url{https://doi.org/10.1137/0729008}

\bibitem{pwlandau1975classical}
Landau, L. D., and Lifshitz, E. M. (1975). \emph{The Classical Theory of Fields}. Pergamon Press. ISBN 978-0-08-025072-4.

\bibitem{pwlieb1973proof}
Lieb, E. H., and Ruskai, M. B. (1973). ``Proof of the strong subadditivity of quantum-mechanical entropy.'' \emph{Journal of Mathematical Physics} \textbf{14}, 1938--1941. \url{https://doi.org/10.1063/1.1666274}

\bibitem{pwlim2020hodge}
Lim, L. H. (2020). ``Hodge Laplacians on graphs.'' \emph{SIAM Review} \textbf{62}, 685--715. \url{https://doi.org/10.1137/18M1223101}

\bibitem{pwlothaire1997combinatorics}
Lothaire, M. (1997). \emph{Combinatorics on Words}. Cambridge University Press. \url{https://doi.org/10.1017/CBO9780511566097}

\bibitem{pwmirsky1971dual}
Mirsky, L. (1971). ``A Dual of Dilworth's Decomposition Theorem.'' \emph{The American Mathematical Monthly} \textbf{78}, 876--877. \url{https://doi.org/10.1080/00029890.1971.11992886}

\bibitem{pwmisner1973gravitation}
Misner, C. W., Thorne, K. S., and Wheeler, J. A. (1973). \emph{Gravitation}. W. H. Freeman. ISBN 0-7167-0344-0.

\bibitem{pwmlodinow2018discrete}
Mlodinow, L., and Brun, T. A. (2018). ``Discrete spacetime, quantum walks, and relativistic wave equations.'' \emph{Physical Review A} \textbf{97}, 042131. \url{https://doi.org/10.1103/PhysRevA.97.042131}

\bibitem{pwmoise1977geometric}
Moise, E. E. (1977). \emph{Geometric Topology in Dimensions 2 and 3}. Springer. \url{https://doi.org/10.1007/978-1-4612-9906-6}

\bibitem{pwmoore2009introduction}
Moore, R. E., Kearfott, R. B., and Cloud, M. J. (2009). \emph{Introduction to Interval Analysis}. SIAM. \url{https://doi.org/10.1137/1.9780898717716}

\bibitem{pwmotzkin1952approximation}
Motzkin, T. S., and Wasow, W. (1952). ``On the approximation of linear elliptic differential equations by difference equations with positive coefficients.'' \emph{Journal of Mathematics and Physics} \textbf{31}, 253--259. \url{https://doi.org/10.1002/sapm1952311253}

\bibitem{pwnaber2012geometry}
Naber, G. L. (2012). \emph{The Geometry of Minkowski Spacetime: An Introduction to the Mathematics of the Special Theory of Relativity}. Springer. \url{https://doi.org/10.1007/978-1-4419-7838-7}

\bibitem{pwnarasimhan2007geometric}
Narasimhan, G., and Smid, M. (2007). \emph{Geometric Spanner Networks}. Cambridge University Press. \url{https://doi.org/10.1017/CBO9780511546884}

\bibitem{pwniederreiter1992random}
Niederreiter, H. (1992). \emph{Random Number Generation and Quasi-Monte Carlo Methods}. SIAM. \url{https://doi.org/10.1137/1.9781611970081}

\bibitem{pwnielsen2010quantum}
Nielsen, M. A., and Chuang, I. L. (2010). \emph{Quantum Computation and Quantum Information: 10th Anniversary Edition}. Cambridge University Press. \url{https://doi.org/10.1017/CBO9780511976667}

\bibitem{pwnzongani2024dirac}
Nzongani, U., Eon, N., M{\'{a}}rquez-Mart{\'{i}}n, I., P{\'{e}}rez, A., Di Molfetta, G., and Arrighi, P. (2024). ``Dirac quantum walk on tetrahedra.'' \emph{Physical Review A} \textbf{110}, 042418. \url{https://doi.org/10.1103/PhysRevA.110.042418}

\bibitem{pwoneill1983semi}
O'Neill, B. (1983). \emph{Semi-Riemannian Geometry with Applications to Relativity}. Academic Press, New York. ISBN 0-12-526740-1.

\bibitem{pwpathria2011statistical}
Pathria, R. K., and Beale, P. D. (2011). \emph{Statistical Mechanics}. Academic Press. ISBN 978-0-12-382188-1.

\bibitem{pwpazy1983semigroups}
Pazy, A. (1983). \emph{Semigroups of Linear Operators and Applications to Partial Differential Equations}. Springer. \url{https://doi.org/10.1007/978-1-4612-5561-1}

\bibitem{pwpenrose1984spinors}
Penrose, R., and Rindler, W. (1984). \emph{Spinors and Space-Time, Volume 1}. Cambridge University Press. \url{https://doi.org/10.1017/CBO9780511564048}

\bibitem{pwpoletto1999linear}
Poletto, M., and Sarkar, V. (1999). ``Linear scan register allocation.'' \emph{ACM Transactions on Programming Languages and Systems} \textbf{21}, 895--913. \url{https://doi.org/10.1145/330249.330250}

\bibitem{pwrado1925begriff}
Rad{\'{o}}, T. (1925). ``{\"{U}}ber den Begriff der Riemannschen Fl{\"{a}}che.'' \emph{Acta Litterarum ac Scientiarum Regiae Universitatis Hungaricae Francisco-Josephinae, Sectio Scientiarum Mathematicarum} \textbf{2}, 101--121. \url{http://acta.bibl.u-szeged.hu/13312/}

\bibitem{pwranard2022converse}
Ranard, D., Walter, M., and Witteveen, F. (2022). ``A Converse to Lieb--Robinson Bounds in One Dimension Using Index Theory.'' \emph{Annales Henri Poincar{\'{e}}} \textbf{23}, 3905--3979. \url{https://doi.org/10.1007/s00023-022-01193-x}

\bibitem{pwratcliffe2019foundations}
Ratcliffe, J. G. (2019). \emph{Foundations of Hyperbolic Manifolds}. Springer International Publishing. \url{https://doi.org/10.1007/978-3-030-31597-9}

\bibitem{pwrenes2004symmetric}
Renes, J. M., Blume-Kohout, R., Scott, A. J., and Caves, C. M. (2004). ``Symmetric informationally complete quantum measurements.'' \emph{Journal of Mathematical Physics} \textbf{45}, 2171--2180. \url{https://doi.org/10.1063/1.1737053}

\bibitem{pwrothenberg1971identification}
Rothenberg, T. J. (1971). ``Identification in Parametric Models.'' \emph{Econometrica} \textbf{39}, 577--591. \url{https://doi.org/10.2307/1913267}

\bibitem{pwsantalo2004integral}
Santal{\'{o}}, L. A. (2004). \emph{Integral Geometry and Geometric Probability}. Cambridge University Press. \url{https://doi.org/10.1017/CBO9780511617331}

\bibitem{pwserafini2017quantum}
Serafini, A. (2017). \emph{Quantum Continuous Variables: A Primer of Theoretical Methods}. CRC Press. \url{https://doi.org/10.1201/9781315118727}

\bibitem{pwserre1977linear}
Serre, J. P. (1977). \emph{Linear Representations of Finite Groups}. Springer. \url{https://doi.org/10.1007/978-1-4684-9458-7}

\bibitem{pwshannon1948mathematical}
Shannon, C. E. (1948). ``A Mathematical Theory of Communication.'' \emph{Bell System Technical Journal} \textbf{27}, 379--423. \url{https://doi.org/10.1002/j.1538-7305.1948.tb01338.x}

\bibitem{pwsharafutdinov1994integral}
Sharafutdinov, V. A. (1994). \emph{Integral Geometry of Tensor Fields}. De Gruyter. \url{https://doi.org/10.1515/9783110900095}

\bibitem{pwsloan1994lattice}
Sloan, I. H., and Joe, S. (1994). \emph{Lattice Methods for Multiple Integration}. Clarendon Press, Oxford. ISBN 0-19-853472-8.

\bibitem{pwsontag1998mathematical}
Sontag, E. D. (1998). \emph{Mathematical Control Theory}. Springer. \url{https://doi.org/10.1007/978-1-4612-0577-7}

\bibitem{pwsparso2001principles}
Spars{\o}, J., and Furber, S. (2001). \emph{Principles of Asynchronous Circuit Design}. Springer. \url{https://doi.org/10.1007/978-1-4757-3385-3}

\bibitem{pwstanley1999enumerative}
Stanley, R. P. (1999). \emph{Enumerative Combinatorics, Volume 2}. Cambridge University Press. \url{https://doi.org/10.1017/CBO9780511609589}

\bibitem{pwstanley2015catalan}
Stanley, R. P. (2015). \emph{Catalan Numbers}. Cambridge University Press. \url{https://doi.org/10.1017/CBO9781139871495}

\bibitem{pwstewart1990matrix}
Stewart, G. W., and Sun, J. g. (1990). \emph{Matrix Perturbation Theory}. Academic Press. ISBN 0-12-670230-6.

\bibitem{pwstrauss1989nonlinear}
Strauss, W. A. (1989). \emph{Nonlinear Wave Equations}. American Mathematical Society (CBMS Regional Conference Series in Mathematics 73). \url{https://doi.org/10.1090/cbms/073}

\bibitem{pwstrikwerda2004finite}
Strikwerda, J. C. (2004). \emph{Finite Difference Schemes and Partial Differential Equations}. SIAM. \url{https://doi.org/10.1137/1.9780898717938}

\bibitem{pwtaflove2005computational}
Taflove, A., and Hagness, S. C. (2005). \emph{Computational Electrodynamics: The Finite-Difference Time-Domain Method}. Artech House. ISBN 1-58053-832-0.

\bibitem{pwtakesaki1970tomita}
Takesaki, M. (1970). \emph{Tomita's Theory of Modular Hilbert Algebras and its Applications}. Springer. \url{https://doi.org/10.1007/BFb0065832}

\bibitem{pwterras1999fourier}
Terras, A. (1999). \emph{Fourier Analysis on Finite Groups and Applications}. Cambridge University Press. \url{https://doi.org/10.1017/CBO9780511626265}

\bibitem{pwverlet1967computer}
Verlet, L. (1967). ``Computer "Experiments" on Classical Fluids. I. Thermodynamical Properties of Lennard-Jones Molecules.'' \emph{Physical Review} \textbf{159}, 98--103. \url{https://doi.org/10.1103/PhysRev.159.98}

\bibitem{pwvillasenor1992rigorous}
Villasenor, J., and Buneman, O. (1992). ``Rigorous charge conservation for local electromagnetic field solvers.'' \emph{Computer Physics Communications} \textbf{69}, 306--316. \url{https://doi.org/10.1016/0010-4655(92)90169-Y}

\bibitem{pwwald1984general}
Wald, R. M. (1984). \emph{General Relativity}. University of Chicago Press. \url{https://doi.org/10.7208/chicago/9780226870373.001.0001}

\bibitem{pwwatrous2018theory}
Watrous, J. (2018). \emph{The Theory of Quantum Information}. Cambridge University Press. \url{https://doi.org/10.1017/9781316848142}

\bibitem{pwweedbrook2012gaussian}
Weedbrook, C., Pirandola, S., Garc{\'{i}}a-Patr{\'{o}}n, R., Cerf, N. J., Ralph, T. C., Shapiro, J. H., et al. (2012). ``Gaussian quantum information.'' \emph{Reviews of Modern Physics} \textbf{84}, 621--669. \url{https://doi.org/10.1103/RevModPhys.84.621}

\bibitem{pwwiesbrock1993half}
Wiesbrock, H. W. (1993). ``Half-sided modular inclusions of von-Neumann-algebras.'' \emph{Communications in Mathematical Physics} \textbf{157}, 83--92. \url{https://doi.org/10.1007/BF02098019}

\bibitem{pwwiesbrock1998modular}
Wiesbrock, H. W. (1998). ``Modular intersections of von Neumann algebras in quantum field theory.'' \emph{Communications in Mathematical Physics} \textbf{193}, 269--285. \url{https://doi.org/10.1007/s002200050329}

\bibitem{pwyee1966numerical}
Yee, K. S. (1966). ``Numerical solution of initial boundary value problems involving Maxwell's equations in isotropic media.'' \emph{IEEE Transactions on Antennas and Propagation} \textbf{14}, 302--307. \url{https://doi.org/10.1109/TAP.1966.1138693}

\bibitem{pwyoshida1990construction}
Yoshida, H. (1990). ``Construction of higher order symplectic integrators.'' \emph{Physics Letters A} \textbf{150}, 262--268. \url{https://doi.org/10.1016/0375-9601(90)90092-3}

\bibitem{pwzeeman1964causality}
Zeeman, E. C. (1964). ``Causality implies the Lorentz group.'' \emph{Journal of Mathematical Physics} \textbf{5}, 490--493. \url{https://doi.org/10.1063/1.1704140}

